\SPBGContentSlide
  {Мягкое отделение: парные проверки}
  {%
    \centering
    \TinyBody
    \renewcommand{\arraystretch}{1.12}
    \setlength{\tabcolsep}{2.8pt}
    \begin{tabular}{|l|r|r|c|r|r|c|}
      \hline
      \textbf{Метод} & \textbf{F1} & \textbf{$\Delta F_1$} & \textbf{$p_H$}
      & \textbf{$t$, c} & \textbf{$\Delta t$} & \textbf{$p_H$} \\
      \hline
      MDM-G & 0.647 & +0.000 & 1.000 & 0.6601 & +0.3605 & $8.77\cdot10^{-7}$ \\
      \hline
      MDM-A & 0.615 & -0.032 & 0.473 & 1.2505 & +0.9509 & $8.77\cdot10^{-7}$ \\
      \hline
      SMO & 0.808 & +0.161 & 0.002 & 0.0362 & -0.2635 & $8.77\cdot10^{-7}$ \\
      \hline
      SVC & 0.808 & +0.161 & 0.002 & 0.0018 & -0.2978 & $8.77\cdot10^{-7}$ \\
      \hline
      LinearSVC & 0.751 & +0.105 & 0.003 & 0.0004 & -0.2992 & $8.77\cdot10^{-7}$ \\
      \hline
    \end{tabular}

    \vspace{.18cm}
    \begin{minipage}{.90\linewidth}
      \SmallBody
      \raggedright
      Сравнение проводится с базовым мягким MDM на N1--N4. MDM-G и MDM-A не
      дают значимого улучшения качества, зато увеличивают время. SVM-решатели
      значимо выигрывают по F1; по ROC-AUC картина такая же для SMO, SVC и
      LinearSVC.
    \end{minipage}
  }
